% Source: https://artofproblemsolving.com/community/c3h1759220_linear_function_problem_middle_level
% Stable identity: aops-topic-1759220-q01
From $2=1+a$ we get $a=1$, so the line is $y=x+1$ and $b=4$. Reflect $A$ in the $x$-axis to $A'(1,-2)$. For every $P$ on the $x$-axis, $AP=A'P$, so $AP+BP=A'P+PB\ge A'B$, with equality exactly when $P$ lies on segment $A'B$.

\begin{asy}
usepackage("amsmath");
size(5.5cm);
defaultpen(fontsize(10pt));
pair A=(1,2),B=(3,4),A1=(1,-2),P=(5/3,0);
draw((-0.5,0)--(4.5,0),Arrow(4));draw((0,-2.6)--(0,5),Arrow(4));
draw((-0.5,0.5)--(4,5),linewidth(0.8));
draw(A--P--B,red+linewidth(0.9));
draw(A1--P,red+dashed);draw(A--A1,gray+dotted);
dot(A);dot(B);dot(A1);dot(P);
label("$A$",A,1.5*W);label("$B(3,4)$",B,SE);label("$A'(1,-2)$",A1,E);label("$P$",P,2*dir(-60));
label("$y=x+1$",(4,5),E);label("$x$",(4.5,0),S);label("$y$",(0,5),W);
label("\mbox{}",(min(currentpicture,user=true)+max(currentpicture,user=true))/2,black); // keeps the text after the figure black when LaTeX typesets the labels inline
\end{asy}

Line $A'B$ has slope $\dfrac{4-(-2)}{3-1}=3$, so it is $y=3(x-1)-2$, which meets $y=0$ at $x=\dfrac{5}{3}$. Hence $P=\boxed{\left(\dfrac{5}{3},0\right)}$, and the minimum of $AP+BP$ is $A'B=\sqrt{4+36}=2\sqrt{10}$.
